\setcounter{section}{15}






  \zwischenueberschrift{Übungsaufgaben}

\inputaufgabe
{}
{





Es sei $R$ ein
\definitionsverweis {kommutativer Ring}{}{.}
Zeige die Äquivalenz folgender Aussagen.
\aufzaehlungzwei { $R$ hat genau ein
\definitionsverweis {maximales Ideal}{}{.}
} {Die Menge der
\definitionsverweis {Nichteinheiten}{}{}
$R \setminus R^\times$ bildet ein
\definitionsverweis {Ideal}{}{}
in $R$.
}





}
{} {}

\inputaufgabe
{}
{





Es sei $R$ ein
\definitionsverweis {kommutativer Ring}{}{.}
Zeige, dass $R$ genau dann ein
\definitionsverweis {lokaler Ring}{}{}
ist, wenn
\mathl{a+b}{} nur dann eine
\definitionsverweis {Einheit}{}{}
ist, wenn $a$ oder $b$ eine Einheit ist.





}
{} {}

\inputaufgabe
{}
{





Bestimme die
\definitionsverweis {Unterringe}{}{}
der
\definitionsverweis {rationalen Zahlen}{}{}
$\Q$, die
\definitionsverweis {lokal}{}{}
sind.





}
{} {}

\inputaufgabe
{}
{





Es sei $R$ ein \definitionsverweis {kommutativer}{}{}
\definitionsverweis {lokaler Ring}{}{.}
Zeige, dass $R$
\definitionsverweis {zusammenhängend}{}{}
ist.





}
{} {}

\inputaufgabe
{}
{





Es sei ${\mathfrak n}$ ein
\definitionsverweis {maximales Ideal}{}{}
in einem
\definitionsverweis {kommutativen Ring}{}{}
$R$. Es sei $R_{\mathfrak n}$ die
\definitionsverweis {Lokalisierung}{}{}
von $R$ an ${\mathfrak n}$ und es sei

\mavergleichskette
{\vergleichskette
{ {\mathfrak m}
}
{ = }{ {\mathfrak n} R_{\mathfrak n}
}
{  }{
}
{    }{
}
{   }{
}
}
{}{}{}
das maximale Ideal von $R_{\mathfrak n}$. Zeige

\mavergleichskette
{\vergleichskette
{ R/ {\mathfrak n}
}
{ = }{ R_{\mathfrak n} /{\mathfrak m}
}
{  }{
}
{    }{
}
{   }{
}
}
{}{}{.}





}
{} {}

\inputaufgabegibtloesung
{}
{





Es sei $R$ ein
\definitionsverweis {kommutativer Ring}{}{}
und sei ${\mathfrak p}$ ein
\definitionsverweis {Primideal}{}{.}
Dann ist der
\definitionsverweis {Restklassenring}{}{}

\mavergleichskette
{\vergleichskette
{S
}
{ = }{ R/{\mathfrak p}
}
{  }{
}
{    }{
}
{   }{
}
}
{}{}{}
ein
\definitionsverweis {Integritätsbereich}{}{}
mit
\definitionsverweis {Quotientenkörper}{}{}

\mavergleichskette
{\vergleichskette
{Q
}
{ = }{ Q(S)
}
{  }{
}
{    }{
}
{   }{
}
}
{}{}{}
und $R_{\mathfrak p}$ ist ein
\definitionsverweis {lokaler Ring}{}{}
mit dem
\definitionsverweis {maximalen Ideal}{}{}

\mathl{{\mathfrak p}R_{\mathfrak p}}{.} Zeige, dass eine natürliche Isomorphie

\mavergleichskettedisp
{\vergleichskette
{ Q(S)
}
{ \cong} { R_{\mathfrak p}/ {\mathfrak p} R_{\mathfrak p}
}
{ } {
}
{ } {
}
{ } {
}
}
{}{}{}
vorliegt.





}
{(Man nennt diesen Körper auch den
\definitionswortenp{Restekörper}{} zu ${\mathfrak p}$).} {}

\inputaufgabe
{}
{





Zeige, dass zu

\mavergleichskette
{\vergleichskette
{ a
}
{ \in }{ {\mathbb C}
}
{  }{
}
{    }{
}
{   }{
}
}
{}{}{}
der
\definitionsverweis {Einsetzungshomomorphismus}{}{}
\maabbeledisp {} { {\mathbb C}[X] } { {\mathbb C}
} { X } {a
} {,}
mit der
\definitionsverweis {Evaluationsabbildung}{}{}
\zusatzklammer {in den
\definitionsverweis {Restekörper}{}{}

\mathdisp {{\mathbb C}[X]_{(X-a)}/ (X-a) {\mathbb C}[X]_{(X-a)}} {}
} {} {}
zum
\definitionsverweis {Primideal}{}{}

\mathl{(X-a)}{} übereinstimmt.





}
{} {}

\inputaufgabe
{}
{





Es sei $R$ ein
\definitionsverweis {lokaler Ring}{}{}
mit
\definitionsverweis {Restekörper}{}{}
$K$. Zeige, dass $R$ und $K$ genau dann die gleiche
\definitionsverweis {Charakteristik}{}{}
haben, wenn $R$ einen
\definitionsverweis {Körper}{}{}
enthält.





}
{} {}

\inputaufgabegibtloesung
{}
{





Es sei $R$ ein
\definitionsverweis {lokaler Ring}{}{}
und ${\mathfrak a}$ ein
\definitionsverweis {Ideal}{}{}
von $R$. Zeige, dass
\maabbdisp {} { R ^{\times} } { \left(  R/ {\mathfrak a}   \right)^{\times}
} {}
\definitionsverweis {surjektiv}{}{}
ist.





}
{} {}

\inputaufgabe
{}
{





Es sei $K$ ein
\definitionsverweis {algebraisch abgeschlossener Körper}{}{}
und

\mavergleichskette
{\vergleichskette
{ R
}
{ = }{ K[X_1 , \ldots , X_n]
}
{  }{
}
{    }{
}
{   }{
}
}
{}{}{.}
Zeige, dass sämtliche
\definitionsverweis {Lokalisierungen}{}{}
von $R$ an
\definitionsverweis {maximalen Idealen}{}{}
zueinander
\definitionsverweis {isomorph}{}{}
sind.





}
{} {}

\inputaufgabe
{}
{





Es sei $K$ ein
\definitionsverweis {Körper}{}{}
und betrachte das Achsenkreuz

\mavergleichskettedisp
{\vergleichskette
{V
}
{ =} { K\!-\!\operatorname{Spek}\, \left(  K[X,Y]/(XY)  \right)
}
{ } {
}
{ } {
}
{ } {
}
}
{}{}{.}
Bestimme für jeden Punkt

\mavergleichskette
{\vergleichskette
{P
}
{ \in }{ V
}
{  }{
}
{    }{
}
{   }{
}
}
{}{}{,}
ob der
\definitionsverweis {lokale Ring}{}{}
an $P$ ein
\definitionsverweis {Integritätsbereich}{}{}
ist oder nicht.





}
{} {}

\inputaufgabe
{}
{





Wir betrachten die Neilsche Parabel

\mavergleichskettedisp
{\vergleichskette
{ C
}
{ =} { V \left(  X^2-Y^3  \right)
}
{ \subseteq} { {\mathbb A}^{2}_{ K }
}
{ } {
}
{ } {
}
}
{}{}{}
über einem
\definitionsverweis {algebraisch abgeschlossenen Körper}{}{}
$K$. Zeige, dass sämtliche
\definitionsverweis {Lokalisierungen}{}{}
von $C$ an Punkten

\mavergleichskette
{\vergleichskette
{ P
}
{ \neq }{ (0,0)
}
{  }{
}
{    }{
}
{   }{
}
}
{}{}{}
zueinander
\definitionsverweis {isomorph}{}{}
sind, aber nicht zur Lokalisierung im Nullpunkt.





}
{} {}

\inputaufgabe
{}
{





Es sei $R$ die
\definitionsverweis {Lokalisierung}{}{}
im Nullpunkt der Kurve

\mavergleichskettedisp
{\vergleichskette
{ C
}
{ =} { V \left(  Y^2-X^2-X^3  \right)
}
{ \subseteq} { {\mathbb A}^{2}_{ K }
}
{ } {
}
{ } {
}
}
{}{}{}
und es sei $S$ die Lokalisierung des Achsenkreuzes im Nullpunkt. Sind diese beiden
\definitionsverweis {lokalen Ringe}{}{}
\definitionsverweis {isomorph}{}{?}





}
{} {}

\inputaufgabe
{}
{





Es sei $R$ ein
\definitionsverweis {kommutativer Ring}{}{}
und sei ${\mathfrak p}$ ein
\definitionsverweis {Primideal}{}{.}
Zeige, dass ${\mathfrak p}$ genau dann ein
\definitionsverweis {minimales Primideal}{}{}
ist, wenn die
\definitionsverweis {Reduktion}{}{}
der
\definitionsverweis {Lokalisierung}{}{}
$R_{\mathfrak p}$ ein
\definitionsverweis {Körper}{}{}
ist.





}
{} {}

\inputaufgabe
{}
{





Es sei $R$ ein kommutativer Ring und sei ${\mathfrak m}$ ein
\definitionsverweis {maximales Ideal}{}{} mit
\definitionsverweis {Lokalisierung}{}{} $R_{\mathfrak m}$. Es sei ${\mathfrak a}$ ein Ideal, das unter der Lokalisierungsabbildung zum Kern gehört. Zeige, dass dann $R_{\mathfrak m}$ auch eine Lokalisierung von $R/{\mathfrak a}$ ist.





}
{} {}

\inputaufgabe
{}
{





Es sei $K$ ein
\definitionsverweis {Körper}{}{} und $R$ eine
\definitionsverweis {endlich erzeugte}{}{}
$K$-\definitionsverweis {Algebra}{}{.}
Es sei

\mavergleichskette
{\vergleichskette
{ S
}
{ = }{ R_{\mathfrak m}
}
{  }{
}
{    }{
}
{   }{
}
}
{}{}{}
die
\definitionsverweis {Lokalisierung}{}{}
von $R$ an einem
\definitionsverweis {maximalen Ideal}{}{}
${\mathfrak m}$. Zeige, dass der
\definitionsverweis {Restekörper}{}{}
von $S$
\definitionsverweis {endlich}{}{}
über $K$ ist.





}
{} {}

\inputaufgabe
{}
{





Es sei $R$ ein
\definitionsverweis {kommutativer Ring}{}{,}
sei

\mavergleichskette
{\vergleichskette
{f
}
{ \in }{ R
}
{  }{
}
{    }{
}
{   }{
}
}
{}{}{}
und sei ${\mathfrak a}$ ein
\definitionsverweis {Ideal}{}{.}
Zeige, dass

\mavergleichskette
{\vergleichskette
{f
}
{ \in }{ {\mathfrak a}
}
{  }{
}
{    }{
}
{   }{
}
}
{}{}{}
genau dann gilt, wenn für alle
\definitionsverweis {Lokalisierungen}{}{}
$R_{\mathfrak p}$ gilt, dass

\mavergleichskette
{\vergleichskette
{f
}
{ \in }{ {\mathfrak a} R_{\mathfrak p}
}
{  }{
}
{    }{
}
{   }{
}
}
{}{}{}
ist.





}
{Bemerkung: Man sagt daher, dass die Idealzugehörigkeit eine lokale Eigenschaft ist.} {}

\inputaufgabe
{}
{





Es sei $R$ ein
\definitionsverweis {kommutativer Ring}{}{.}
Beweise die Äquivalenz folgender Aussagen.
\aufzaehlungdrei{$R$ ist
\definitionsverweis {reduziert}{}{.}
}{Für jedes
\definitionsverweis {Primideal}{}{}

\mathl{{\mathfrak p}}{} ist $R_{\mathfrak p}$ reduziert.
}{Für jedes
\definitionsverweis {maximale Ideal}{}{}

\mathl{{\mathfrak m}}{} ist $R_{\mathfrak m}$ reduziert.
}





}
{Bemerkung: Man sagt daher, dass die Reduziertheit eine lokale Eigenschaft ist.

Man gebe auch ein Beispiel für einen kommutativen Ring, der nicht integer ist, dessen Lokalisierungen an Primidealen aber alle integer sind.} {}

\inputaufgabegibtloesung
{}
{





Es sei $K$ ein
\definitionsverweis {Körper}{}{}
und seien $R$ und $S$ integre,
\definitionsverweis {endlich erzeugte}{}{}
$K$-\definitionsverweis {Algebren}{}{.}
Es sei
\maabbdisp {\varphi} { R } { S
} {}
ein
$K$-\definitionsverweis {Algebrahomomorphismus}{}{}
und ${\mathfrak n}$ ein
\definitionsverweis {maximales Ideal}{}{}
in $S$ mit

\mavergleichskette
{\vergleichskette
{ \varphi^{-1}( {\mathfrak n})
}
{ = }{ {\mathfrak m}
}
{  }{
}
{    }{
}
{   }{
}
}
{}{}{.}
Die Abbildung induziere einen
\definitionsverweis {Isomorphismus}{}{}
\maabb {} { R_{\mathfrak m} } { S_{\mathfrak n}
} {.}
Zeige, dass es dann auch ein

\mathbed {f \in R} {}
 {f \not \in {\mathfrak m}} {}
 {} {} {} {,}
derart gibt, dass
\maabb {} { R_f } { S_{\varphi (f)}
} {}
ein Isomorphismus ist.





}
{} {}

\inputaufgabe
{}
{





Es sei $K$ ein
\definitionsverweis {algebraisch abgeschlossener Körper}{}{}
und $R$ eine
\definitionsverweis {endlich erzeugte}{}{}
kommutative
$K$-\definitionsverweis {Algebra}{}{.}
Es seien $F_1$ und $F_2$
\definitionsverweis {topologische Filter}{}{}
in
\mathl{K\!-\!\operatorname{Spek}\, \left(  R  \right)}{} mit

\mavergleichskette
{\vergleichskette
{ F_1
}
{ \subseteq }{ F_2
}
{  }{
}
{    }{
}
{   }{
}
}
{}{}{.}
Zeige, dass es einen
\definitionsverweis {Ringhomomorphismus}{}{}
\maabbdisp {} { {\mathcal O}_{F_1 } } { {\mathcal O}_{F_2 }
} {}
gibt.





}
{} {}

\inputaufgabe
{}
{





Es sei $K$ ein
\definitionsverweis {algebraisch abgeschlossener Körper}{}{}
und $R$ eine
\definitionsverweis {endlich erzeugte}{}{}
kommutative $K$-Algebra. Es sei

\mavergleichskette
{\vergleichskette
{ P
}
{ \in }{ K\!-\!\operatorname{Spek}\, \left(  R  \right)
}
{  }{
}
{    }{
}
{   }{
}
}
{}{}{}
ein Punkt. Zeige
\zusatzklammer {ohne

Satz 15.12

zu verwenden} {} {}, dass der
\definitionsverweis {Halm}{}{}
${\mathcal O}_{ P }$ ein
\definitionsverweis {lokaler Ring}{}{}
ist.





}
{} {}

\inputaufgabegibtloesung
{}
{





Es sei $K$ ein Körper und $R$ eine integre, endlich erzeugte $K$-Algebra mit Quotientenkörper
\mathl{Q(R)}{.} Es sei $q \in Q(R)$. Zeige, dass die Menge

\mathdisp {\left\{  P \in K\!-\!\operatorname{Spek}\, \left(  R  \right)   \mid   q \in {\mathcal O}_P         \right\}} {  }

offen in $K\!-\!\operatorname{Spek}\, \left(  R  \right)$ ist
\zusatzklammer {dabei bezeichnet ${\mathcal O}_P$ den lokalen Ring im Punkt $P$} {} {.}





}
{} {}

\inputaufgabe
{}
{





Es sei $I$ eine
\definitionsverweis {gerichtete Indexmenge}{}{}
und sei $G_i$,
\mathl{i \in I}{,} ein
\definitionsverweis {gerichtetes System}{}{}
von kommutativen Gruppen. Zeige, dass der
\definitionsverweis {Kolimes}{}{}
eine kommutative Gruppe ist.





}
{} {}

\inputaufgabe
{}
{





Es sei $I$ eine
\definitionsverweis {gerichtete Indexmenge}{}{}
und sei $M_i$,
\mathl{i \in I}{,} ein
\definitionsverweis {gerichtetes System}{}{}
von Mengen. Es sei $N$ eine weitere Menge und zu jedem

\mavergleichskette
{\vergleichskette
{ i
}
{ \in }{ I
}
{  }{
}
{    }{
}
{   }{
}
}
{}{}{}
sei eine Abbildung
\maabbdisp {\psi_i} { M_i } { N
} {}
mit der Eigenschaft gegeben, dass

\mavergleichskette
{\vergleichskette
{ \psi_i
}
{ = }{ \psi_j \circ \varphi_{ij}
}
{  }{
}
{    }{
}
{   }{
}
}
{}{}{}
ist für alle

\mavergleichskette
{\vergleichskette
{ i
}
{ \preccurlyeq }{ j
}
{  }{
}
{    }{
}
{   }{
}
}
{}{}{}
\zusatzklammer {wobei $\varphi_{ij}$ die Abbildungen des Systems bezeichnen} {} {.}
Beweise die universelle Eigenschaft des
\definitionsverweis {Kolimes}{}{,}
nämlich, dass es eine eindeutig bestimmte Abbildung
\maabbdisp {\psi} { \operatorname{colim}\,_{i \in I} M_i  } { N
} {}
derart gibt, dass

\mavergleichskette
{\vergleichskette
{ \psi_i
}
{ = }{ \psi  \circ j_i
}
{  }{
}
{    }{
}
{   }{
}
}
{}{}{}
ist, wobei
\mathl{j_i:M_i \rightarrow \operatorname{colim}\,_{i \in I} M_i}{} die natürlichen Abbildungen sind.

Zeige ferner, dass falls $M_i$ ein gerichtetes System von Gruppen und falls $N$ ebenfalls eine Gruppe ist und alle $\psi_i$ Gruppenhomomorphismen sind, dass dann auch $\psi$ ein Gruppenhomomorphismus ist.





}
{} {}





  \zwischenueberschrift{Aufgaben zum Abgeben}

\inputaufgabe
{4}
{





Beschreibe die Menge $M$ aller
\mathl{2 \times 3}{-}Matrizen mit Rang $\leq 1$ über einem Körper $K$ als $K$-Spektrum einer geeigneten $K$-Algebra. Zeige, dass es eine Isomorphie zwischen einer (nicht leeren) Zariski-offenen Teilmenge von $M$ und einer offenen Menge des ${ {\mathbb A}_{ K }^{  4  } }$ gibt.





}
{} {}

\inputaufgabe
{4}
{





Es sei $K$ ein
\definitionsverweis {Körper}{}{}
und sei $R$ eine
$K$-\definitionsverweis {Algebra von endlichem Typ}{}{.}
Es seien
\mathl{P_1 , \ldots , P_n}{} endlich viele Punkte in

\mavergleichskette
{\vergleichskette
{ X
}
{ = }{ K\!-\!\operatorname{Spek}\, \left(  R  \right)
}
{  }{
}
{    }{
}
{   }{
}
}
{}{}{}
Zeige, dass der
\definitionsverweis {Umgebungsfilter}{}{}
dieser Punkte durch offene Mengen der Form $D(f)$ erzeugt wird.





}
{D.h. es ist zu zeigen, dass es zu

\mavergleichskettek
{\vergleichskettek
{ P_1 , \ldots ,  P_n
}
{ \in }{ U
}
{  }{
}
{    }{
}
{   }{
}
}
{}{}{}
offen stets ein

\mavergleichskette
{\vergleichskette
{ F
}
{ \in }{ R
}
{  }{
}
{    }{
}
{   }{
}
}
{}{}{}
mit

\mavergleichskettek
{\vergleichskettek
{ P_1 , \ldots ,  P_n
}
{ \in }{  D(F)
}
{ \subseteq }{ U
}
{    }{
}
{   }{
}
}
{}{}{}
gibt.} {}

\inputaufgabe
{5 (1+2+2)}
{





Es sei $K$ ein
\definitionsverweis {Körper}{}{,}
sei $R$ eine kommutative
$K$-\definitionsverweis {Algebra von endlichem Typ}{}{}
und sei $S$ ein
\definitionsverweis {multiplikatives System}{}{}
in $R$. Zu $S$ definieren wir

\mavergleichskettedisp
{\vergleichskette
{ F(S)
}
{ =} { \left\{  U \subseteq   K\!-\!\operatorname{Spek}\, \left(  R  \right)   \text{ offen}  \mid   \text{ es gibt } f \in S \text{ mit } D(f) \subseteq U         \right\}
}
{ } {
}
{ } {
}
{ } {
}
}
{}{}{.}
\aufzaehlungdreiabc{Zeige, dass

\mavergleichskette
{\vergleichskette
{ F
}
{ = }{ F(S)
}
{  }{
}
{    }{
}
{   }{
}
}
{}{}{}
ein
\definitionsverweis {topologischer Filter}{}{}
in
\mathl{K\!-\!\operatorname{Spek}\, \left(  R  \right)}{} ist.
}{Zeige, dass es einen
\definitionsverweis {Ringhomomorphismus}{}{}
\maabbdisp {} { R_S  } { { \mathcal O}_F
} {}
gibt.
}{Zeige, dass der Ringhomomorphismus aus (b) eine
\definitionsverweis {Isomorphie}{}{}
ist, falls $K$
\definitionsverweis {algebraisch abgeschlossen}{}{}
und $R$
\definitionsverweis {reduziert}{}{}
ist.
}





}
{} {}

\inputaufgabe
{4}
{





Es sei

\mavergleichskette
{\vergleichskette
{ X
}
{ = }{ K\!-\!\operatorname{Spek}\, \left(  R  \right)
}
{  }{
}
{    }{
}
{   }{
}
}
{}{}{}
eine
\definitionsverweis {affine Varietät}{}{}
und seien

\mavergleichskette
{\vergleichskette
{ P_1 , \ldots , P_n
}
{ \in }{ X
}
{  }{
}
{    }{
}
{   }{
}
}
{}{}{}
endlich viele Punkte. Es sei $F$ der
\definitionsverweis {Umgebungsfilter}{}{}
dieser Punkte und ${\mathcal O}_F$ der zugehörige
\definitionsverweis {Halm}{}{.}
Zeige, dass ${\mathcal O}_F$ genau dann ein lokaler Ring ist, wenn

\mavergleichskette
{\vergleichskette
{ n
}
{ = }{ 1
}
{  }{
}
{    }{
}
{   }{
}
}
{}{}{}
ist.





}
{} {}

\inputaufgabe
{4}
{





Es sei $R$ ein
\definitionsverweis {kommutativer Ring}{}{}
und sei

\mavergleichskette
{\vergleichskette
{S
}
{ \subseteq }{ R
}
{  }{
}
{    }{
}
{   }{
}
}
{}{}{}
ein
\definitionsverweis {multiplikatives System}{}{.}
Auf $S$ betrachten wir folgende
\zusatzklammer {partielle} {} {}
\definitionsverweis {Ordnung}{}{,}
und zwar sagen wir

\mavergleichskette
{\vergleichskette
{f
}
{ \preccurlyeq }{g
}
{  }{
}
{    }{
}
{   }{
}
}
{}{}{,}
falls $f$ eine Potenz von $g$ teilt
\zusatzklammer {und wir identifizieren zwei Elemente, wenn diese Relation in beide Richtungen vorliegt} {} {.}
Zeige, dass die kommutativen Ringe

\mathdisp {R_f, \, f \in S} { , }

ein
\definitionsverweis {gerichtetes System}{}{}
bilden, und dass für den
\definitionsverweis {Kolimes}{}{}

\mavergleichskettedisp
{\vergleichskette
{ \operatorname{colim}_{f \in S} R_f
}
{ =} { R_S
}
{ } {
}
{ } {
}
{ } {
}
}
{}{}{}
gilt.





}
{} {}
